\paragraph{Proof of Theorem~\ref{jointdiam:main}.}
The definitions give $0<\rho_-\le\rho_+$; continuity of $V$ ensures
that the maximum exists. For a nonoptimal source,
write $r=\|x-x^*\|$ and let $p$ be its terminal radial derivative.
Shifted RSI and constrained minimality give
$mr\le p\le P$, while $pr\le G(x)\le g$ and $r\le D$. Hence
\[
 pr\le\min\{g,P^2/m,PD\}=h,\qquad
 r\le R:=\min\{\sqrt{h/m},D\}.
\]
Put $r_0=h/P\le R$. The radial envelope $A(P,r)$ in
\eqref{joint:radial-identity} is strictly increasing in $r$.
For $r\le r_0$ it bounds the loss by
$A(P,r_0)=V(h,r_0^2)$. For $r\ge r_0$, the terminal cap $h/r$
and initial cap $h/r-mr$ bound it by $V(h,r^2)$. This proves
\eqref{jointdiam:definition}. A zero gap or zero norm forces zero
error by shifted RSI. The zero bound is attained at the minimizer
of a quadratic on a triangle of diameter $D$.

For sharpness, choose a maximizing radius $\rho$, and put
\[
 r=\sqrt\rho,\quad p=h/r,\quad q=p-mr,\quad
 c=\sqrt{P^2-p^2},\quad a_D=\sqrt{D^2-r^2/4}.
\]
The scalar construction in Lemma~\ref{compact:radial} gives
smooth functions $u_j$ on neighborhoods of $[0,r]$ with
\[
 \begin{gathered}
 u_j(0)=0,\quad u_j'(0)=q,\quad u_j'(r)=p,\quad
 u_j'(s)-q\ge ms,\\
 |u_j''(s)|\le H_0+H_1u_j(s),\qquad
 u_j(r)\longrightarrow V(h,\rho).
 \end{gathered}
\]
Define
\[
 Q=\operatorname{conv}\{(0,0),(r,0),(r/2,a_D)\},\qquad
 f_j(s,y)=u_j(s)+cy+my^2/2.
\]
The two sloping sides have length $D$ and the base has length
$r\le D$, so $Q$ has diameter exactly $D$ and nonempty interior.
Its points satisfy $0\le s\le r$ and $y\ge0$. The unique minimizer
is $(0,0)$, and throughout $Q$,
\[
 \langle\nabla f_j(s,y)-\nabla f_j(0,0),(s,y)\rangle
       =(u_j'(s)-q)s+my^2\ge m(s^2+y^2),
\]
\[
 \|\nabla^2f_j(s,y)\|
       =\max\{|u_j''(s)|,m\}
       \le H_0+H_1u_j(s)\le H_0+H_1f_j(s,y).
\]
At $(r,0)$ the gradient is $(p,c)$, of norm $P$. Its linear
pairing is minimized at $(0,0)$, so the full gap is $pr=h\le g$.
The source loss tends to $\mathcal J_D(g,P)$. This proves sharpness
for upper bounds, even with diameter fixed exactly.

For exact positive observations, $g\le PD$ is necessary by
$G(x)\le\|\nabla f(x)\|\operatorname{diam}Q$. If also
$g\le P^2/m$, then $h=g$, and the preceding construction already
has both observations exact. If $P^2/m<g\le PD$, set
\[
 r=P/m,\qquad a=g/P-r,\qquad \ell=r+a=g/P\le D.
\]
Use smooth scalar profiles on $[0,r]$ with initial slope zero
and terminal slope $P$, approaching $\mathcal E(P)$. Their
smoothing is confined to the interior switch, so near zero
they agree with
\[
 u_+(s)=
 \begin{cases}
 (H_0/H_1)\{\cosh(\sqrt{H_1}s)-1\},&H_1>0,\\
 H_0s^2/2,&H_1=0.
 \end{cases}
\]
Extend them to $[-a,0]$ by this even function. There
$|u_+''|=H_0+H_1u_+$ and $u_+'(s)s\ge ms^2$.
Denote the extension by $u_j$.
On the triangle
\[
 Q=\operatorname{conv}\left\{(-a,0),(r,0),
       \left((r-a)/2,\sqrt{D^2-\ell^2/4}\right)\right\}
\]
put $f_j(s,y)=u_j(s)+my^2/2$. Again $\operatorname{diam}Q=D$,
the minimizer is $(0,0)$, and shifted RSI and the Hessian bound
hold on all of $Q$. The source gradient is $(P,0)$; the linear
oracle selects $(-a,0)$, giving gap $P\ell=g$. The loss tends
to $\mathcal E(P)=\mathcal J_D(g,P)$. All the functions are
smooth on neighborhoods of their triangles.

Formula \eqref{jointdiam:constant} maximizes the affine radius
profile $h+\beta\rho$. Formula \eqref{jointdiam:ordinary} follows
from Corollary~\ref{gapregime:radius-monotone} at $\rho_-$.
For Corollary~\ref{jointdiam:norm}, the gap bound $PD$ follows
from the norm and diameter bounds. In
$\mathcal J_D(PD,P)$ the radius interval is the singleton
$r_D^2$, and its initial slope is
$P-mr_D=\max\{P-mD,0\}=q$. The radial identity and
Proposition~\ref{joint:fixed-slope} give
\eqref{jointdiam:norm-formula}; the constant-curvature radius
formula gives \eqref{jointdiam:norm-constant}.
